\begingroup\fontsize{9.25}{11.4}\selectfont
\setlength{\tabcolsep}{3pt}\renewcommand{\arraystretch}{1.18}
\begin{longtable}{@{}>{\raggedright\arraybackslash}p{0.38\linewidth}rp{0.16\linewidth}p{0.26\linewidth}@{}}
\caption{Project composition and assignment jumps through 2012}\label{tab:pub_A04} \\
\toprule Project group & Records / CCPP & Share (\%) & Assignment jump (SE), pp \\ \midrule\endfirsthead
\multicolumn{4}{c}{\tablename\ \thetable{} -- continued} \\
\toprule Project group & Records / CCPP & Share (\%) & Assignment jump (SE), pp \\ \midrule\endhead
Productive and livelihood & 106 / 106 & 53.27 & 51.51 (25.84) \\
Social and basic services & 41 / 41 & 20.60 & 19.00 (18.83) \\
Community and civic infrastructure & 35 / 35 & 17.59 & 6.92 (3.70) \\
Management and capacity support & 17 / 17 & 8.54 & 0.68 (0.85) \\
\bottomrule\end{longtable}
\parbox{0.97\linewidth}{\footnotesize \textit{Notes:} Composition is descriptive among 199 CMAN project records linked to the 487-community SISFOH analysis universe and recorded by 2012. Communities may appear in more than one project group. Assignment jumps are robust bias-corrected local-linear discontinuities in receiving at least one project of the indicated group, in percentage points, using \(h=0.0075\), triangular weights, mass-point adjustment, and district CR2 inference. Project type is a post-assignment implementation attribute; treated-only outcome comparisons are not labeled causal. Sources: RUV, CMAN, and SISFOH 2012--2013.}
\par\smallskip{\footnotesize\textit{Review boundary:} Conditional local evidence; historical design search and selection remain limitations. RF denotes the assignment reduced form; CCPP denotes a community. In outcome tables, shares, binary outcomes and zero-to-one indices are in percentage points; logarithms remain log points and counts remain counts. A common window is not claimed optimal for every outcome. No artifact-level release approval or manuscript synchronization is implied.}\par\endgroup
